\chapter{User Testing and Further Improvements}\label{ch:user_testing_improvements}
This chapter presents the user testing survey, the parts of the system that were fixed
based on the testing feedback and possible future improvements.


\section{User Testing}
After the platform was implemented, a short user testing survey was
performed. The goal was to see how real users interact with the
system, whether they are able to easily navigate and use the available features,
and which parts of the interface are difficult to understand.

\subsection{Methods}
The survey was again created in Google Forms. It contained 10 tasks,
which were based on the existing use cases. Together, these tasks
covered the main features available to a typical user, such as
registration, music playback, and social actions. The exact wording
of the tasks is available in the survey form.

After each task, the participant selected one of three outcomes:
``Completed without help'', ``Completed but needed help or hints'', or
``Failed / gave up''. The participant also rated the difficulty of the
task from 1 (very easy) to 5 (very hard). An optional text field was
included after each task, so participants could add comments.

The form is available
online,\footnote{\url{https://docs.google.com/forms/d/19hc0Pa-td0c6V1uaFZowmoxGE1ID0-wfhmEqBcgC0CY/}}
and the raw answers are stored in \path{text/testing_survey_results.csv}.

\subsection{Results}
In total, 11 people took part in the survey. Most tasks were completed
without help, and the average difficulty was close to 1.4 on the
1 to 5 scale. This shows that the main user flows were generally easy
to use.

However, not all parts of the system were equally clear. Some tasks
were less intuitive, and several comments pointed on concrete usability
issues. These issues are described in the next subsection.

\subsection{Discovered Issues}
The list below describes the main issues found during testing. It also
explains whether each issue was fixed after the survey.

\begin{itemize}
    \item Creating a new Post was not intuitive, because the action was
    available only from the User profile, and the Posts menu itself did not
    include a clear way to create a new Post. This was improved by
    adding a button directly to the Posts menu. In addition, the User's
    own Posts are now displayed there as well, which makes the flow
    easier to discover.

    \item Icon-only buttons across the app had no hover labels. Because
    of this, participants had to click the buttons to understand what
    they did. Tooltips were added to the header, player bar, Song and
    Collection menus, Notifications, and the search window.

    \item Email normalization was inconsistent. The Backend stored only the
    domain part in lowercase, while the login endpoint compared the whole email in
    lowercase. This made emails such as \path{User@example.com}
    succeed in registration but fail at login. The
    Backend now normalizes the whole address.

    \item Toast pop-ups appeared at the bottom of the page and overlapped
    the playback controls. They were moved to the top middle part of the
    screen.

    \item Two issues were left for future work. First, the Queue exit
    mechanic was not obvious, because the Queue is currently opened and
    closed by the same navigation button. Second, the volume slider and
    audio normalization did not work in Safari.

\end{itemize}

Overall, the testing showed that the main flows work well, but it also
revealed several usability gaps. Most of these issues were fixed after
testing. The remaining problems, together with other possible future
work, are described in the following section.


\section{Further Improvements}\label{sec:improvements}
This section lists known gaps and possible next steps for the
app.

\subsection{Out-of-Scope Items}\label{subsec:imp_out_of_scope}
On the Frontend, several features are missing. There is no
mobile layout and no offline Playback, and the interface is
English-only. Accessibility is also only partial. For example,
some clickable elements use \texttt{div} instead of
\texttt{button}, and the custom controls have no Accessible Rich
Internet Applications~(ARIA) roles or labels. There is also no
Digital Rights Management~(DRM) mechanism. In addition,
Web Audio API (audio normalization, volume controls) does not work in Safari.\footnote{Even though \texttt{caniuse} lists Safari as a compatible browser for Web Audio API\parencite{caniuse_webaudio}}\parencite{webkit_hls_webaudio,shaka_safari_webaudio}

Some smaller issues also remain. The Queue view could use a clearer
close or back action. In Safari, audio normalization and volume controls
based on the Web Audio API did not work reliably during
testing.\footnote{Even though \texttt{caniuse} lists Safari as a compatible
browser for Web Audio API\parencite{caniuse_webaudio}}\parencite{webkit_hls_webaudio,shaka_safari_webaudio}

Moreover, a few Backend features have no Frontend
implementation. For example, there is no UI for the password
reset and password change endpoints.
Similarly, the Queue view does not support drag-to-reorder,
even though the Backend has a reorder endpoint.

Finally, production troubleshooting relies on reading the
server logs over SSH\@. There is no error reporter, no metrics
pipeline, and no structured log format.

\subsection{Playback Synchronization}\label{subsec:imp_clock}
The clock calculation method from
\hyperref[sec:playback]{Playback and Device Handling} only takes
the difference between the client and server clocks at receive
time. As a result, on slow networks the estimate is off by about
the one-way packet trip time. A better method would record both
the send time and the receive time of each broadcast, use half
the round-trip as the offset, and average several samples (this
is the approach NTP uses).~\parencite{rfc5905_ntp}

\subsection{Backend Refactors}\label{subsec:imp_be_refactors}
The three Notification managers send
\texttt{notifications.changed} inside
\texttt{Manager.create}. As a result, the data layer is coupled
to the WebSocket layer. The broadcast call should move into the
view or a small service object, the same way other parts of the
app do it.

A similar problem appears in the
\hyperref[sec:publications]{Publications} section, which uses
one Django model for Posts, Collection Comments, Chat messages,
and Replies. The data is similar, but the behavior for each
kind is very different. For example, Collection Comments are
only allowed on public non-draft Collections, and Posts are only
allowed on the User's own Feed. Chat messages, in turn, need
\texttt{Chat} and \texttt{ChatMember} support. The API is
already split into one endpoint per kind. Therefore, the next
step is to split the model in the same way Notifications were
split.

\subsection{Frontend Cache Strategy}\label{subsec:imp_cache}
Almost every WebSocket event handler on the Frontend calls\\
\texttt{queryClient.invalidateQueries}, even when the event
payload already has the new state. As a result, every change
forces a refetch and shows a short loading state. Where the
payload is enough, the handler should call
\texttt{queryClient.setQueryData} instead, and keep invalidation
only as a fallback. This would remove the extra round-trip.

\subsection{Test Coverage}\label{subsec:imp_tests}
Backend test coverage is not high. The Notifications app has no
tests at all, the audio conversion tests only cover lossless
source inputs, and the endpoints are not tested thoroughly.
Moreover, the Frontend has no automated tests at all. Component
tests with React Testing Library and some smoke tests for the
main flows would be a good next step.
